\section{Branching, with the flow-first construction inside it}
\label{sec:branching}

The branching hierarchy is read off the robustness criterion of
Section~\ref{sec:rmp}. A row with $\phi\equiv1$ --- any bound on a sum of
dispatch counts, including the per-class total $n_{gt}$ and the fleet
$X=\sum_{g,t}n_{gt}$ --- has its dual folded into $\mu_{gt}$ exactly as
\eqref{bp:MP-count} does, and all such aggregates are integer in every
plan. A row with $\phi$ affine in the image, a bound on
$\overline{V}_{nt}$ or $\overline{V^{\mathrm{e}}}_{rt}$, is absorbed
into $\pi_{ct}$, already an input of the pricing problem; these count
physical RTIs, so they too may carry a dichotomy. Bounding a single
$\lambda^{g}_{\omega t}$ is neither: honouring it means excluding that
column from the subproblem
\autocite{Vanderbeck2000OnIntegerProgramming}, destroying the property
that one column set serves every period. It is also useless here, the
master holding many columns that differ in a single RTI, so the
relaxation answers by shifting the same fractional weight one column
over --- measured at $3\,000$ branches with no integral relaxation
reached.

Between the two robust rungs the first binds, because transport cost
is $cK\cdot X$ and nothing else in the objective counts trucks. A bound
on a delivery is answered by moving freight between periods at the
price of one period's holding cost, and so is a bound on a
\emph{single} period's dispatch total; the aggregate with no such
escape is the horizon total. We therefore branch on the fleet, in two
levels (Figure~\ref{fig:tree}), keeping the delivery rung as a third
entered only where the first two have nothing left to say. Master,
pricing problem, construction and tree are together the procedure we
call BP-3D-SCPTOP (Algorithm~\ref{alg:bp3d}).

\paragraph{Level 1: the total fleet.}
Given an incumbent with fleet vector $\mathbf{n}^{\star}$ and total
$X^{\star}$, three regions partition the space. $X\ge X^{\star}+1$
costs a whole truck with no offset and is usually pruned by bound;
$X=X^{\star}$ fixes transport, so everything inside it is inventory
allocation and this is where level 2 goes; and $X\le X^{\star}-1$ is
the only region that can beat the incumbent outright, so it is worth a
certificate of its own and the stack is ordered to expand it first. A
plan found \emph{inside} an inequality region splits it on its own
total the same way, so level 1 recurses.

\paragraph{Level 2: the fleet vector, inside $X=X^{\star}$.}
Fix an order on the pairs $(g,t)$. The fleet vector of any plan in the
region differs from $\mathbf{n}^{\star}$ at a first index $j$, or
nowhere, so branching there gives three children,
\begin{equation}
n_{j}\le n^{\star}_{j}-1,
\qquad
n_{j}= n^{\star}_{j},
\qquad
n_{j}\ge n^{\star}_{j}+1,
\label{eq:murty}
\end{equation}
disjoint and exhaustive, with the equality child partitioned again at
$j+1$ (Murty's scheme \autocite{Murty1968AnAlgorithm}) and expanded
first. The partition is taken around the plan \emph{constructed at that
node}, so a node that improves the incumbent is re-partitioned around
the better plan at once and nothing is restarted.

% F6 of article/paper-figures.md, redrawn, and the only drawing of the
% branching scheme.  One picture for the whole hierarchy: the three rungs,
% nested (level 2 inside the equality child of level 1, level 3 inside the
% Murty leaf); an integer line under the first split, so that "disjoint and
% exhaustive" is shown rather than asserted and the later rungs can be read
% off it; the construction as what happens at every node; and the rung that
% is deliberately not taken.  Greyscale only, straight edges only, every
% label clear of every edge, and the spine of the tree runs at x=7.5 -- the
% heading blocks are 66mm wide so that none of them reaches it.
\begin{figure}[!t]
\centering
\begin{tikzpicture}[>=Latex, x=1cm, y=1cm,
  nd/.style={rectangle, draw=black!75, fill=white, line width=0.5pt,
             inner sep=3pt, minimum width=26mm, minimum height=6mm,
             font=\footnotesize},
  fx/.style={nd, fill=black!12},
  rt/.style={nd, minimum width=52mm},
  rep/.style={rectangle, draw=black!60, densely dashed, fill=white,
              line width=0.5pt, inner sep=3pt, text width=40mm,
              align=flush center, font=\scriptsize},
  an/.style={font=\scriptsize, align=flush center, text width=60mm, anchor=north},
  hd/.style={font=\scriptsize, align=flush left, anchor=north west,
             text width=66mm},
  bnr/.style={rectangle, draw=black!45, fill=black!5, rounded corners=1pt,
              inner sep=4pt, text width=142mm, align=flush left, font=\scriptsize,
              anchor=north},
  bnx/.style={bnr, draw=black!60, densely dashed, fill=white},
  sp/.style={draw=black!75, line width=0.5pt},
  ed/.style={sp, ->},
  ax/.style={draw=black!55, line width=0.4pt},
  seg/.style={draw=black!85, line width=1.5pt},
  tk/.style={font=\tiny, inner sep=1pt}]

% ---------------------------------------------------------------- banner --
\node[bnr] (ban) at (7.5,1.62) {\textbf{At every node} the construction
  (Algorithm~\ref{alg:construct}) builds the plan that node is partitioned
  around, and its trucks join one shared pool: a column is valid in every
  region, so a pruned node leaves its columns behind.};

% =============================================================== level 1 ==
\node[hd] at (0.4,0.62) {\textbf{level 1} --- the total fleet
  $X=\sum_{g,t}n_{gt}$};

\node[rt] (root) at (7.5,-0.45)
  {root: the constructed plan $\mathbf{n}^{\star}$, $X^{\star}$ trucks};

\node[nd] (a1) at (3.0,-1.70) {$X\le X^{\star}-1$};
\node[fx] (b1) at (7.5,-1.70) {$X=X^{\star}$};
\node[nd] (c1) at (12.0,-1.70) {$X\ge X^{\star}+1$};
\coordinate (j1) at (7.5,-1.07);
\draw[sp] (root.south) -- (j1);
\draw[ed] (j1) -- (a1.north);
\draw[ed] (j1) -- (b1.north);
\draw[ed] (j1) -- (c1.north);

% the integer line: the three children abut at $X^\star$ and leave no gap
\draw[ax] (0.95,-2.55) -- (14.05,-2.55);
\foreach \i in {0,...,14} {\draw[ax] (1.2+\i*0.9,-2.55) -- ++(0,0.075);}
\draw[seg] (1.2,-2.55) -- (6.6,-2.55);
\draw[seg] (8.4,-2.55) -- (13.8,-2.55);
\fill (7.5,-2.55) circle (1.7pt);
\node[tk, anchor=east] at (1.05,-2.55) {$\cdots$};
\node[tk, anchor=west] at (13.95,-2.55) {$\cdots$};
\node[tk] at (6.6,-2.82) {$X^{\star}\!-\!1$};
\node[tk] at (8.4,-2.82) {$X^{\star}\!+\!1$};
\node[tk, anchor=west] at (7.66,-2.42) {$X^{\star}$};
\node[tk, anchor=east] at (13.8,-2.82) {$X\in\mathbb{Z}$};

\node[an] at (3.0,-3.07) {the only region that can hold a cheaper plan:
  popped first, and split again on its own total};
\node[an] at (12.0,-3.07) {a whole truck more, with nothing in the objective
  to offset it: usually pruned by bound};

% =============================================================== level 2 ==
\node[hd] at (0.4,-4.00) {\textbf{level 2} --- the fleet vector, inside
  $X=X^{\star}$: Murty at the first pair $j$ that may differ};

\node[nd] (a2) at (3.0,-5.45) {$n_{j}\le n^{\star}_{j}-1$};
\node[fx] (b2) at (7.5,-5.45) {$n_{j}=n^{\star}_{j}$};
\node[nd] (c2) at (12.0,-5.45) {$n_{j}\ge n^{\star}_{j}+1$};
\coordinate (j2) at (7.5,-4.85);
\draw[sp] (b1.south) -- (7.5,-2.55);
\draw[sp] (7.5,-2.55) -- (j2);
\draw[ed] (j2) -- (a2.north);
\draw[ed] (j2) -- (b2.north);
\draw[ed] (j2) -- (c2.north);

\node[an, text width=38mm] at (3.0,-5.95) {the same trichotomy, one pair at
  a time; empty when $n^{\star}_{j}=0$};
\node[an, text width=38mm] at (12.0,-5.95) {the total is fixed here, so a pair
  that moves up forces another down};

\node[rep] (rep) at (7.5,-6.37) {split again at the next pair,\\
  until every $n_{gt}$ is fixed};
\draw[ed] (b2.south) -- (rep.north);

% =============================================================== level 3 ==
\node[hd] at (0.4,-7.20) {\textbf{level 3} --- one delivery, once every
  $n_{gt}$ is fixed: the one furthest from an integer};

\node[nd, minimum width=32mm] (a3) at (4.9,-8.65)
  {$\overline{V}_{nt}\le\lfloor v\rfloor$};
\node[nd, minimum width=32mm] (c3) at (10.1,-8.65)
  {$\overline{V}_{nt}\ge\lceil v\rceil$};
\coordinate (j3) at (7.5,-8.05);
\draw[sp] (rep.south) -- (j3);
\draw[ed] (j3) -- (a3.north);
\draw[ed] (j3) -- (c3.north);

\node[an, text width=104mm] at (7.5,-9.25) {a delivery counts physical RTIs,
  so the two children are disjoint and cover $\mathbb{Z}$ between them; this
  rung is entered only where the first two have nothing left to say};

% --------------------------------------------------------- the non-rung --
\node[bnx] (ban2) at (7.5,-10.19) {\textbf{Not a level:} a single column
  variable $\lambda^{g}_{\omega t}$. A bound on one is not a row of the
  master --- honouring it means excluding that column from the pricing
  problem --- and the relaxation answers by moving the same weight to a
  column that differs in one RTI: $3\,000$ branches, no integral relaxation
  reached.};

\end{tikzpicture}
\caption{Every rung bounds an integer quantity the master already prices,
so the children of a split are disjoint and together exhaustive and the search
gives up nothing by taking them. The order is forced by the objective:
transport costs $cK\cdot X$ and nothing else in the model counts trucks, so
the horizon total is the one aggregate that cannot be answered by moving
freight between periods for one period's holding cost, as a single dispatch
count or a single delivery can. And because these are rows of the master and
not restrictions on the subproblem, branching moves the duals and leaves the
pricing problem alone --- which is also how a load no dual vector would have
asked for (Figure~\ref{fig:hull}) comes to be priced at all.}
\label{fig:tree}
\end{figure}


\paragraph{Pruning, and what counts as a proof.}
A node is a region imposed on the RMP as fleet rows. Its LP over the
current pool is solved first as a filter (Section~\ref{sec:bounds});
where the filter allows pruning and the certification policy admits the
node, column generation is run to a certificate, and a certified value
at or above the incumbent prunes it while a value below becomes that
region's bound. Every other node carries its parent's, valid because
the child's region is a subset. An infeasible node relaxation is
\emph{not} a proof that a region is empty: stock non-negativity is not
elastic, so a fleet lower bound can force deliveries up through
whatever columns the pool holds and kill the relaxation for a reason
internal to the pool. Only arithmetic licenses that sentence, and three
such arguments are available: the fleet box intersected with the
interval on $X$, an availability cap telescoped from the sending rows
\eqref{sb:eq:inve-plant} and \eqref{sb:eq:inv-sup-main} under
non-negative stock, and the constants of \eqref{eq:covercut} read as a
requirement floor.

\paragraph{What the tree reports.}
The incumbent is the cheapest verified plan any construction found. The
global bound is the minimum over the \emph{frontier}, not the root, a
region nobody expanded contributing nothing; at the deadline the open
regions are swept and certified where the filter allows. Under 3D
loading a region's certificate comes, like the root's, from the
volumetric relaxation with the region's fleet rows applied, which
Proposition~\ref{prop:inclusion} covers because those rows bound the
fleet vector and not the load model. A column, by contrast, is a
feasible load, so one found in any region is valid in every other and a
pruned node leaves its columns behind. Optimality is claimed
\emph{only} when the root is certified, no region remains open, no
restart has broken the partition and the incumbent meets the global
bound. The hierarchy is finite and its levels exhaustive, so those
conditions are reachable in principle; no run here meets them.

\begin{algorithm}[!htb]
\caption{BP-3D-SCPTOP: branch and bound over the fleet vector, with
\textsc{Construct} as primal heuristic and \textsc{Certify} as bound.}
\label{alg:bp3d}
\begin{algorithmic}[1]
\State $\bar\Omega\gets$ one maximally full single-commodity truck per commodity
   and direction, each verified
   \Comment{no demand mixes, no reduced variants}
\State $\mathrm{LB}_{0}\gets\textsc{Certify}(\bar\Omega,\text{root})$;
   $(\bar z,\mathbf{n}^{\star},\bar\Omega)\gets\textsc{Construct}(\bar\Omega,\text{root})$;
   $X^{\star}\gets\sum_{g,t}n^{\star}_{gt}$
\State $\mathcal{S}\gets[\,X\ge X^{\star}+1,\; X=X^{\star},\; X\le X^{\star}-1\,]$
   \Comment{a stack: the region that can beat the incumbent is popped first}
\While{$\mathcal{S}\ne\emptyset$ and time remains}
   \State $\mathcal{R}\gets$ pop $\mathcal{S}$; impose it on the RMP as fleet rows
   \If{$\mathcal{R}$ is refuted arithmetically} \State \textbf{continue}
      \Comment{fleet box, availability, requirement} \EndIf
   \State $\tilde z\gets$ LP over $\bar\Omega$ inside $\mathcal{R}$
      \Comment{over-estimates the region's LP: a filter, never a bound}
   \If{$\tilde z\ge\bar z$ and the policy certifies $\mathcal{R}$}
      \State $(\mathrm{LB}_{\mathcal{R}},\text{proved})\gets
             \textsc{Certify}(\bar\Omega,\mathcal{R})$
      \If{$\mathrm{LB}_{\mathcal{R}}\ge\bar z$} \State \textbf{continue}
         \Comment{pruned} \EndIf
      \State record $\mathrm{LB}_{\mathcal{R}}$ as an open bound
   \Else \State $\mathcal{R}$ carries its parent's bound \EndIf
   \State $(z,\hat{\mathbf{n}},\bar\Omega)\gets\textsc{Construct}(\bar\Omega,\mathcal{R})$
      \Comment{its columns join the shared pool whatever it returns}
   \If{no plan} \State price under $\mathcal{R}$'s duals and construct once
      more; \textbf{continue} if still none \EndIf
   \If{$z<\bar z$} \State $(\bar z,\mathbf{n}^{\star})\gets(z,\hat{\mathbf{n}})$
      \Comment{new incumbent} \EndIf
   \If{$\mathcal{R}$ does not fix $X$}
      \State push the three regions of level 1 around
            $\sum_{g,t}\hat n_{gt}$ \Comment{level 1 recurses}
   \ElsIf{the Murty depth $j$ is below the cap}
      \State push the three children \eqref{eq:murty} at the $j$-th pair,
             around $\hat{\mathbf{n}}$, equality last
   \EndIf
\EndWhile
\ForAll{$\mathcal{R}$ still on $\mathcal{S}$, while a reserve of time remains}
   \State filter by the pool LP; certify the survivors; record their bounds
   \Comment{the frontier}
\EndFor
\State $\mathrm{LB}\gets\max\{\mathrm{LB}_{0},\,
   \min\text{ over the open and frontier bounds}\}$, capped at $\bar z$
\State verify every dispatched truck of the incumbent against
   \eqref{eq:col-geom}--\eqref{eq:col-agg}
\State \Return $(\bar z,\mathrm{LB})$; optimality is claimed only if
   $\mathrm{LB}_{0}$ was proved, $\mathcal{S}=\emptyset$, no restart occurred
   and $\bar z=\mathrm{LB}$
\end{algorithmic}
\end{algorithm}


\subsection{The initial pool}
\label{sec:initcols}
The search starts from the plainest possible pool: one maximally full
truck per commodity and direction, obtained geometrically by pricing
that commodity alone with every other capped at zero --- so that where
the probe closes, the pool truck is that commodity's placeable maximum
and the same number is its placeability cap --- and volumetrically by
filling to capacity. No demand mixes and no reduced variants are
seeded, so whatever the search reaches it reached by pricing, branching
and construction.

\subsection{The flow-first construction}
\label{sec:relaxfix-flow}

The primal heuristic decides first how much moves and how many trucks
carry it, and builds the trucks afterwards. It runs inside the region
$\mathcal{R}$ of the node that calls it and must stay there, or the
children that node induces would partition around a point it excludes,
and it runs on a \emph{fresh copy} of the restricted master, step~2
changing variable types the node's own LP needs intact.
Figure~\ref{fig:construct} follows one period through the four steps of
Algorithm~\ref{alg:construct}.

% F5 of article/paper-figures.md: the flow-first construction, one period,
% top to bottom.  The object passing between the steps is drawn, not named,
% and every band carries the same two-line statement of what it fixes and
% what it leaves free.  Greyscale only (fills black!X plus one hatch), so
% nothing depends on colour.  The libraries used -- patterns,
% decorations.pathreplacing, arrows.meta -- are all loaded by
% ../necessarypackages and paper.tex; no new package is needed.
\begin{figure}[!htb]
\centering
\resizebox{\textwidth}{!}{%
\begin{tikzpicture}[>=Latex, line width=0.5pt,
  band/.style={draw=black!45, fill=black!2, rounded corners=2pt},
  ttl/.style={font=\small\bfseries, anchor=west, align=left},
  bdg/.style={circle, draw=black!70, fill=white, inner sep=0pt,
              minimum size=4.6mm, font=\scriptsize\bfseries},
  ann/.style={font=\scriptsize, align=center},
  annl/.style={font=\scriptsize, align=left},
  mic/.style={font=\tiny, align=center},
  micl/.style={font=\tiny, align=left},
  st/.style={font=\scriptsize, align=left, anchor=west, text width=38mm},
  barO/.style={draw=black!70, fill=white},         % still continuous
  qty/.style={draw=black!70, fill=black!25},
  trk/.style={draw=black!80, fill=white, rounded corners=0.6pt},
  seg1/.style={draw=black!55, fill=black!30},      % R1 units
  seg2/.style={draw=black!55, fill=black!10},      % R2 units
  sur/.style={draw=black!65, pattern=north east lines, pattern color=black!75},
  glyA/.style={draw=black!55, fill=black!12},      % a column already in the pool
  glyB/.style={draw=black!70, fill=black!35},      % a part built here
  rule/.style={draw=black!35, dotted}]

\def\BW{15.2}          % content width
\def\SX{10.95}         % the fixed/free rule
\def\GX{2.20}          % x of the between-band arrows

% =====================================================================
% band 1 -- the node's LP
% =====================================================================
\draw[band] (0,0) rectangle (\BW,-2.25);
\draw[rule] (\SX,-0.12) -- (\SX,-2.13);
\node[bdg] at (0.55,-0.38) {1};
\node[ttl] at (0.95,-0.38) {Restricted master, inside $\mathcal{R}$};

\draw[black!55] (0.60,-1.82) -- (4.55,-1.82);
\foreach \i/\h/\v in {0/0.95/{1.0}, 1/0.86/{0.9}, 2/0.67/{0.7},
                      3/0.48/{0.5}, 4/0.29/{0.3}} {
  \pgfmathsetmacro{\x}{0.85+\i*0.72}
  \draw[barO] (\x,-1.82) rectangle (\x+0.46,-1.82+\h);
  \node[mic, anchor=south] at (\x+0.23,-1.79+\h) {\v};
}
\node[mic] at (2.58,-2.00) {one bar per column $\omega\in\bar\Omega$};
\node[annl, anchor=west] at (4.80,-1.38)
  {$\sum_{\omega}\lambda^{g}_{\omega t}=3.4$};
\node[mic, anchor=north west] at (4.80,-1.62) {not an integer};

\node[ann] at (7.55,-0.98) {duals $\pi_{ct},\,\mu_{gt}$};
\draw[->] (7.55,-1.20) -- (7.55,-1.76);
\node[mic, anchor=north] at (7.55,-1.80) {to the pricing problem};

\node[ann, text width=22mm] at (9.60,-1.26)
  {on a \emph{fresh copy} of the RMP: step~2 changes variable types};

\node[st] at (11.15,-1.12)
  {\textbf{fixed:} the region $\mathcal{R}$ --- bounds on $n_{gt}$ and $X$\\[2pt]
   \textbf{free:} every $\lambda^{g}_{\omega t}$, $\overline{V}_{nt}$,
   $\overline{V^{\mathrm{e}}}_{rt}$, all continuous};

% ---- 1 -> 2 ----------------------------------------------------------
\draw[->, line width=0.8pt] (\GX,-2.32) -- (\GX,-2.80);
\node[micl, anchor=west] at (\GX+0.22,-2.56)
  {the same rows; only the variable types change};

% =====================================================================
% band 2 -- the aggregate integer master
% =====================================================================
\draw[band] (0,-2.87) rectangle (\BW,-5.42);
\draw[rule] (\SX,-2.99) -- (\SX,-5.30);
\node[bdg] at (0.55,-3.25) {2};
\node[ttl] at (0.95,-3.25) {Aggregate integer master};

\draw[black!55] (0.60,-5.00) -- (4.55,-5.00);
\foreach \i/\h/\v in {0/0.90/{1.2}, 1/0.68/{0.9}, 2/0.60/{0.8},
                      3/0.53/{0.7}, 4/0.30/{0.4}} {
  \pgfmathsetmacro{\x}{0.85+\i*0.72}
  \draw[barO] (\x,-5.00) rectangle (\x+0.46,-5.00+\h);
  \node[mic, anchor=south] at (\x+0.23,-4.97+\h) {\v};
}
\node[ann, anchor=south] at (2.58,-3.92)
  {$\sum_{\omega}\lambda^{g}_{\omega t}=n^{\star}_{gt}=4$};
\node[mic] at (2.58,-5.18) {the $\lambda$ themselves stay continuous};

\node[ann, anchor=south] at (6.60,-3.92) {$Q^{\star}_{gt}$, integer};
\node[mic, anchor=east] at (5.30,-4.33) {$R_{1}$};
\draw[qty] (5.38,-4.46) rectangle (7.38,-4.20);
\node[mic, anchor=west] at (7.46,-4.33) {200};
\node[mic, anchor=east] at (5.30,-4.83) {$R_{2}$};
\draw[qty] (5.38,-4.96) rectangle (5.68,-4.70);
\node[mic, anchor=west] at (5.76,-4.83) {30};
\node[mic] at (6.60,-5.18) {bar length $\propto$ units};

\node[ann, text width=25mm] at (9.48,-4.42)
  {\emph{no pricing call here}: the rows are the master's own, and the
   objective is decided now};

\node[st] at (11.15,-4.15)
  {\textbf{fixed:} $n^{\star}_{gt}$, and the integer deliveries
   $Q^{\star}_{gt}=(\overline{V}_{nt},\overline{V^{\mathrm{e}}}_{rt})$\\[2pt]
   \textbf{free:} the columns $\lambda^{g}_{\omega t}$, \emph{still} continuous};

% ---- 2 -> 3 ----------------------------------------------------------
\draw[->, line width=0.8pt] (\GX,-5.49) -- (\GX,-5.97);
\node[micl, anchor=west] at (\GX+0.22,-5.73)
  {$(Q^{\star}_{gt},\,n^{\star}_{gt})$ --- fixed over continuous $\lambda$,
   so \emph{not} a sum of whole legal loads};

% =====================================================================
% band 3 -- the partition, as a cover
% =====================================================================
\draw[band] (0,-6.04) rectangle (\BW,-11.46);
\draw[rule] (\SX,-6.16) -- (\SX,-8.98);
\node[bdg] at (0.55,-6.42) {3};
\node[ttl] at (0.95,-6.42) {Partition --- a cover, not a split};

% what has to be delivered
\draw[seg1] (0.72,-8.52) rectangle (1.32,-7.07);
\draw[seg2] (0.72,-7.07) rectangle (1.32,-6.77);
\node[mic] at (1.02,-7.80) {$R_{1}$};
\node[mic, anchor=west] at (1.40,-6.92) {$R_{2}$};
\node[mic, anchor=north] at (1.02,-8.58) {$Q^{\star}_{gt}$};
\draw[->] (1.95,-7.80) -- (2.62,-7.80);
\node[mic, anchor=east] at (2.56,-7.20) {$\alpha\Gamma_{K}$};
\node[mic, anchor=north] at (2.28,-7.88) {$n^{\star}_{gt}=4$};

% the four parts
\foreach \k/\f/\r/\s/\z/\b in {%
   1/1.56/1.12/0.14/{$z_{1}{=}1$}/{full by vol.},
   2/1.14/0.86/0.12/{$z_{2}{=}0$}/{full by wgt.},
   3/1.48/1.04/0.22/{$z_{3}{=}1$}/{full by vol.},
   4/1.06/0.78/0.14/{$z_{4}{=}0$}/{full by wgt.}} {
  \pgfmathsetmacro{\x}{2.75+(\k-1)*1.60}
  \pgfmathsetmacro{\yr}{-8.52+\r}
  \pgfmathsetmacro{\yf}{-8.52+\f}
  \pgfmathsetmacro{\ys}{\yf-\s}
  \draw[trk] (\x,-8.52) rectangle (\x+1.30,-6.77);
  \draw[seg1] (\x+0.06,-8.46) rectangle (\x+1.24,\yr);
  \draw[seg2] (\x+0.06,\yr) rectangle (\x+1.24,\ys);
  \draw[sur]  (\x+0.06,\ys) rectangle (\x+1.24,\yf);
  \draw[dashed, black!85] (\x-0.10,-7.20) -- (\x+1.40,-7.20);
  \node[mic, text width=15mm, anchor=north] at (\x+0.65,-8.60) {\z\\ \b};
}

% key
\draw[sur]  (9.14,-7.14) rectangle (9.44,-6.90);
\node[mic, anchor=west] at (9.52,-7.02) {surplus};
\draw[seg2] (9.14,-7.64) rectangle (9.44,-7.40);
\node[mic, anchor=west] at (9.52,-7.52) {$R_{2}$ units};
\draw[seg1] (9.14,-8.14) rectangle (9.44,-7.90);
\node[mic, anchor=west] at (9.52,-8.02) {$R_{1}$ units};

\node[st] at (11.15,-7.50)
  {\textbf{fixed:} $n^{\star}_{gt}$ trucks ($+1$ or $+2$ if none can be
   formed), and $Q^{\star}_{gtc}$ as a \emph{lower} bound\\[2pt]
   \textbf{free:} the counts $a_{kc}\le\bar u_{gc}$, and which resource
   binds ($z_{k}$)};

\node[ann, text width=144mm, anchor=north] at (7.45,-9.06)
  {$\textstyle\sum_{k}a_{kc}\ge Q^{\star}_{gtc}$ \eqref{eq:part-cover}:
   a \textbf{cover}, not a split. Each part is a legal full truck
   \eqref{eq:part-cap}--\eqref{eq:part-floor}, full on volume \textbf{or}
   on payload and never on the sum of the two; the surplus is minimised in
   volume \eqref{eq:part-obj}, and the parts are ordered by weight
   \eqref{eq:part-sym}.};

\node[ann, text width=139mm, draw=black!45, dashed, fill=black!5,
      rounded corners=1.5pt, inner sep=4pt, anchor=north] at (7.60,-10.18)
  {\emph{Why the inequality cannot be made an equality:}
   $\{R_{1}{:}\,200,\,R_{2}{:}\,30\}$ asks for $35.5\,\mathrm{m}^{2}$ of
   stacks on a $32.64\,\mathrm{m}^{2}$ floor, so not one truck, and holds
   $1.83$ full-truckload floors of volume, so not two full ones --- no
   number of full trucks splits it exactly.};

% ---- 3 -> 4 ----------------------------------------------------------
\draw[->, line width=0.8pt] (\GX,-11.53) -- (\GX,-12.01);
\node[micl, anchor=west] at (\GX+0.22,-11.77)
  {$a_{1},\dots,a_{n^{\star}_{gt}}$ --- each a legal full truck, hence a
   legal column};

% =====================================================================
% band 4 -- the integer master over pool and parts
% =====================================================================
\begin{scope}[shift={(0,-0.66)}]
\draw[band] (0,-11.42) rectangle (\BW,-14.30);
\draw[rule] (\SX,-11.54) -- (\SX,-14.18);
\node[bdg] at (0.55,-11.80) {4};
\node[ttl] at (0.95,-11.80) {Integer master over pool and parts};

\draw[draw=black!40, fill=black!3, rounded corners=1.5pt]
   (0.62,-12.34) rectangle (4.86,-13.06);
\foreach \i in {0,...,5} {
  \pgfmathsetmacro{\x}{0.78+\i*0.40}
  \draw[glyA] (\x,-12.90) rectangle (\x+0.32,-12.50);
}
\foreach \i in {6,...,9} {
  \pgfmathsetmacro{\x}{0.78+\i*0.40}
  \draw[glyB] (\x,-12.90) rectangle (\x+0.32,-12.50);
}
\node[ann, anchor=north] at (1.60,-13.12) {pool $\bar\Omega\,\cup\,\{a_{k}\}$};
\draw[decorate, decoration={brace, amplitude=3pt, mirror}, black!60]
   (3.18,-13.14) -- (4.70,-13.14);
\node[mic, anchor=north] at (3.94,-13.24) {the four parts};

\draw[->, line width=0.8pt] (5.02,-12.70) -- (6.52,-12.70);
\node[mic, text width=22mm, anchor=north] at (6.02,-12.90)
  {integer master, warm-started with the partition};

\foreach \k/\f/\l in {1/{glyB}/{$a_{1}$}, 2/{glyB}/{$a_{2}$},
                      3/{glyB}/{$a_{3}$}, 4/{glyA}/{$\omega\in\bar\Omega$}} {
  \pgfmathsetmacro{\x}{6.85+(\k-1)*1.00}
  \draw[\f] (\x,-13.06) rectangle (\x+0.75,-12.34);
  \node[mic] at (\x+0.375,-12.70) {\l};
}
\node[ann, anchor=south] at (8.73,-12.30) {the node's plan};

\draw[draw=black!45, dashed] (7.05,-14.12) rectangle (7.80,-13.62);
\node[mic] at (7.42,-13.87) {$a_{4}$};
\node[annl, text width=26mm, anchor=west] at (7.96,-13.87)
  {built, not dispatched};

\node[st] at (11.15,-12.90)
  {\textbf{fixed:} the pool $\bar\Omega\cup\{a_{k}\}$, and the region
   $\mathcal{R}$\\[2pt]
   \textbf{free:} integer $\lambda^{g}_{\omega t}$ --- how many of each
   column goes out};
\end{scope}

\end{tikzpicture}}
\caption{The flow-first construction over one period, and why its third
step covers rather than splits. Steps~1 and~2 settle how much moves and
how many trucks carry it while the column variables stay
\emph{continuous}, so the delivery vector they fix need not be a sum of
whole legal loads --- for the vector shown, no number of full trucks
realises it exactly. Step~3 therefore asks only that the parts
\emph{cover} it \eqref{eq:part-cover}, paying for the hatched surplus
\eqref{eq:part-obj}, and makes each part full on volume \emph{or} on
payload \eqref{eq:part-floor} rather than on the sum of the two. Nothing
built is binding: the parts enter step~4 as columns with a warm start,
and the integer master may dispatch fewer of them than the cover
produced.}
\label{fig:construct}
\end{figure}


\begin{enumerate}
\item \emph{Restricted master.} The LP over the pool, inside
  $\mathcal{R}$, gives the duals and the fractional point.

\item \emph{Aggregate integer master.} Introduce the dispatch totals
  $n_{gt}$ as explicit integer variables, make
  $\overline{V}_{nt},\overline{V^{\mathrm{e}}}_{rt}$ integer and leave
  the column variables \emph{continuous}, with $\mathcal{R}$ imposed on
  the $n_{gt}$ as well as on the columns and, under 3D loading, the
  stack-count rows \eqref{eq:stackrows} added. The objective is decided
  here, and \emph{without a single call to a pricing problem}: the row
  set is the master's own. Artificial slack above a tolerance means the
  region is infeasible over this pool, and the node is abandoned.

\item \emph{Partition, as a cover.} For each direction and period,
  realise the integer deliveries $Q^{\star}_{gt}$ as $n^{\star}_{gt}$
  legal trucks by an integer programme over per-truck commodity counts
  $a_{kc}$, $k\le n^{\star}_{gt}$:
  \begin{align}
    \min\;\;&\sum_{c}\gamma_{gc}\Bigl(\sum_{k}a_{kc}-Q^{\star}_{gtc}\Bigr)
      \label{eq:part-obj}\\
    \text{s.t.}\;\;&\sum_{k}a_{kc}\;\ge\;Q^{\star}_{gtc} &&\forall c
      \label{eq:part-cover}\\
    &\textstyle\sum_{c}\gamma_{gc}a_{kc}\le\Gamma_{K},\qquad
     \sum_{c}w_{c}a_{kc}\le pK^{w} &&\forall k
      \label{eq:part-cap}\\
    &\textstyle\sum_{c}\gamma_{gc}a_{kc}\ge\alpha\Gamma_{K}\theta_{k},\qquad
     \sum_{c}w_{c}a_{kc}\ge\alpha pK^{w}(1-\theta_{k}) &&\forall k
      \label{eq:part-floor}\\
    &\textstyle\sum_{c}w_{c}a_{kc}\;\ge\;\sum_{c}w_{c}a_{k+1,c}
      &&\forall k<n^{\star}_{gt}
      \label{eq:part-sym}\\
    &a_{kc}\in\{0,\dots,\bar u_{gc}\},\qquad \theta_{k}\in\{0,1\}.\nonumber
  \end{align}
  \eqref{eq:part-cap}--\eqref{eq:part-floor} make each part a legal
  truck, $\theta_{k}$ selecting the resource on which part $k$ meets the
  floor; \eqref{eq:part-sym} orders the parts by weight; $\bar u_{gc}$
  is the placeability cap below. If $n^{\star}_{gt}$ parts cannot be
  formed, $n^{\star}_{gt}+1$ and $+2$ are tried where $\mathcal{R}$
  allows --- and a timeout is not a proof that none can be.

\item \emph{Integer master over pool plus parts.} The parts join the
  pool as columns and the integer master is solved over it,
  warm-started with the partition; its solution is the node's plan.
\end{enumerate}

\paragraph{Why the partition must be a cover.}
This is the step at which the procedure could quietly become wrong.
With integer $\lambda$ the aggregate
$\sum_{\omega}b_{\omega n}\lambda^{p}_{\omega t}$ is a sum of placeable
loads, so an unrealisable delivery vector cannot be requested; step~2
breaks that link by fixing the aggregates over \emph{continuous}
$\lambda$. A fractional blend of legal loads need not be realisable as
legal loads for \emph{any} number of trucks: floor area is consumed in
whole footprints while \eqref{eq:part-floor} is measured in volume, so
splitting a vector into more trucks consumes more floor without adding
volume and splitting it into fewer violates the trailer. On one
benchmark instance $\{R_{1}\!:\!200,R_{2}\!:\!30\}$ holds $1.83$
full-truckload floors of volume but needs $35.5\,\mathrm{m}^{2}$ of
stacks on a $32.64\,\mathrm{m}^{2}$ floor. A cover is sound because
every part is a legal column and the parts enter step~4 \emph{as
columns with a warm start}: the integer master may dispatch fewer than
were built, and over-delivering is feasible, stock having no upper
bound.

\begin{algorithm}[!htb]
\caption{$\textsc{Construct}(\bar\Omega,\mathcal{R})$ --- the flow-first
relax and fix. Lines 9--14 run under 3D loading only. $\varepsilon_{2}$
is step~2's relative MIP gap and \emph{tol} the tolerance above which
artificial slack means the region admits no plan over this pool.}
\label{alg:construct}
\begin{algorithmic}[1]
\State $\bar\Omega'\gets$ a fresh copy of the RMP over $\bar\Omega$
   \Comment{not the node's own LP: step 2 changes variable types}
\State solve $\mathrm{LP}(\bar\Omega',\mathcal{R})$
   \Comment{duals $\pi,\mu$ and the fractional point the node stands on}
\State make $n_{gt},\overline{V}_{nt},\overline{V^{\mathrm{e}}}_{rt}$ integer;
   keep $\lambda$ continuous; impose $\mathcal{R}$ on $n_{gt}$ and on $\lambda$
\State add the aggregate form of \eqref{eq:stackrows}
   \Comment{3D only; layer coefficient $=$ the height count, the crush
      limit being its own row}
\State solve to relative gap $\varepsilon_{2}$
\If{no solution, or artificials $>$ tol}
   \State \Return $\bot$ \Comment{node abandoned; the region is infeasible over this pool} \EndIf
\State $(Q^{\star},n^{\star})\gets$ the integer deliveries and fleet
\ForAll{$(g,t)\in G\times T$}
   \For{$m=n^{\star}_{gt},\,n^{\star}_{gt}+1,\,n^{\star}_{gt}+2$ while
        $\mathcal{R}$ allows}
      \State solve $\textsc{Partition}(Q^{\star}_{gt},m)$,
             \eqref{eq:part-obj}--\eqref{eq:part-sym}, with $\bar u_{gc}$ as
             variable bounds
      \State \textbf{break} if feasible
      \Comment{a timeout is not a proof of infeasibility}
   \EndFor
   \ForAll{parts $a_{k}$ of the cover}
      \State $(\text{verdict},\text{layout})\gets\textsc{Pack}(a_{k})$
      \If{placeable} \State verify the layout;
          $\bar\Omega\gets\bar\Omega\cup\{(a_{k},\text{layout})\}$ \EndIf
      \If{\emph{proved} unplaceable} \State bisect $\sigma\in(0,1]$ and cut
          every load dominating $\lfloor\sigma a_{k}\rfloor$
          \Comment{validity} \EndIf
      \If{undecided} \State exclude $a_{k}$ as a single point
          \Comment{heuristic restriction, never validity} \EndIf
   \EndFor
   \State re-solve the partition if a cut was written and rounds remain
\EndFor
\State $\bar\Omega\gets\bar\Omega\cup\{\text{the parts}\}$; solve the integer
   master over $\bar\Omega$, warm-started with the partition
\State \Return $(z,\hat{\mathbf{n}},\bar\Omega)$
   \Comment{the parts are kept even when no plan is returned}
\end{algorithmic}
\end{algorithm}


\subsection{The construction under 3D loading}
\label{sec:geo3d}

Rows \eqref{eq:part-cap}--\eqref{eq:part-floor} are exactly the
volumetric column conditions, so over that oracle every part is a
column by construction. Under 3D loading a part must also satisfy
\eqref{eq:col-geom}--\eqref{eq:col-agg}; since the volumetric
conditions are \emph{necessary} for geometric feasibility
(Proposition~\ref{prop:inclusion}) the partition programme is kept as a
\emph{relaxation} of the geometric one, and each part is then tested
for placeability. Step~2 makes no oracle call, so the objective is
still decided before any geometry is examined.

Everything here rests on one property. Placeability is \textbf{downward
closed}: given a layout of $a\ge b$, removing the surplus units from
their slots leaves positions unchanged, so nothing overlaps, every
stack is no taller and no heavier, and an emptied slot is dropped. The
full-truckload floor is no part of placeability, so it does not disturb
the closure.

\paragraph{Deciding placeability.}
A part is submitted to an oracle that returns one of three answers. A
constructive \emph{shelf design} --- a small integer programme laying
the part's stacks in shelves along the trailer, one integer count per
shelf, footprint and orientation --- supplies essentially every
positive answer, and its layout passes the independent verifier of
Section~\ref{sec:res-setup} before it counts. The pricing problem, in
feasibility mode with the true slot bound \eqref{sb:eq:Ibound},
supplies the \emph{infeasibility} proofs.

\paragraph{Two kinds of cut, which must never merge.}
A part \emph{proved} unplaceable yields a \emph{dominance cut}
forbidding every load that dominates it componentwise, valid by the
closure above, so one cut removes the whole up-set of the refuted load
--- including every point of the one-unit-at-a-time staircase the
partition would otherwise walk down. Before cutting, the load is
tightened along its own ray: $\lfloor\sigma a\rfloor$ is probed for
$\sigma$ bisected in $(0,1]$ and the cut written at the smallest scaled
point still \emph{proved} unplaceable. A part left \emph{undecided} is
instead excluded as a single point so that the partition proposes
something else: a restriction of a primal heuristic, never a validity
claim. Reading an undecided verdict as a dominance cut is how a method
comes to forbid its own optimum.

\paragraph{Placeability caps.}
The largest number of one commodity a truck can hold is a scalar and,
by the same closure, monotone, so it is computed rather than walked to:
bisected once per direction and commodity at the start of a run, a
proved optimum tightening it and an undecided probe leaving the
analytic cap, then handed to the partition as $\bar u_{gc}$.

\paragraph{A stack relaxation inside the partition.}
The partition additionally carries, per truck $k$ and footprint class
$f$, an integer number of stacks
$s_{kf}\le\lfloor pK\cdot qK/p_{f}q_{f}\rfloor$, the crush limit and
the height budget imposed per stack and the stacks' areas summed
against the floor:
\begin{equation}
\sum_{c\in f}w_{c}a_{kc}\le W^{\max}s_{kf},
\qquad
\sum_{c\in f}r_{gc}\,a_{kc}\le rK\,s_{kf},
\qquad
\sum_{f}p_{f}q_{f}\,s_{kf}\le pK\cdot qK ,
\label{eq:stackrows}
\end{equation}
and step~2 carries the same rows per (direction, period) with the area
right-hand side multiplied by $n_{gt}$. Both are necessary, every truck
of the period having integer stack counts that satisfy the per-truck
rows, and they are what the volumetric relaxation omits. Where a class
carries one RTI type the height row sharpens to
$\sum_{c\in f}a_{kc}\le\lfloor rK/r_{gr}\rfloor s_{kf}$: the crush
limit is already carried by the first row with each commodity's own
weight, so the height row takes the height-limited count, which is what
the runs reported here used; the smaller of the two counts is also
valid, and strictly tighter. These aggregates remain a relaxation of
the per-truck split, which is why step~3 keeps its $n+1$ fallback.